% ============================================================================
% NOVA SECÇÃO para o capítulo ESTADO DA ARTE.
% Colocar ANTES da secção "SÍNTESE E LACUNAS IDENTIFICADAS", que passa a ser
% a última do capítulo.
% ============================================================================

\section{SISTEMAS PROPOSTOS NO DESAFIO SOCCERNET-GSR}
\label{sec:sistemas_gsr}

As secções anteriores percorreram a literatura por tema, reunindo trabalhos publicados
em conferências e revistas. Sucede que a reconstrução do estado de jogo é uma tarefa
recente e que boa parte do trabalho mais relevante não se encontra nesses canais, mas
sim nos relatórios anuais do desafio SoccerNet. São sistemas completos, construídos
sobre o mesmo conjunto de dados, a mesma \emph{baseline} e a mesma métrica, o que faz
deles o termo de comparação mais direto para esta dissertação. Realizaram-se até à data
duas edições da tarefa, documentadas por Cioppa et al. (2024) e por Giancola et al.
(2025), e é sobre elas que incide esta secção.

Antes de as percorrer, convém fixar uma ressalva que se aplica a todos os valores aqui
citados. A classificação oficial é apurada no conjunto de \emph{challenge}, cujas
anotações são reservadas, e não no conjunto de teste, de anotações públicas, que é o
usado neste trabalho. Os valores servem, por isso, para caracterizar abordagens e
tendências, e não para comparação numérica direta.

\subsection{EDIÇÃO DE 2024}

A primeira edição reuniu 16 equipas e 34 submissões, e elevou o GS-HOTA de 23.36\% da
\emph{baseline} para 63.81\% da solução vencedora.

A equipa \textbf{Constructor tech} (63.81\%), cuja abordagem é detalhada em Golovkin et al.
(2025), reconstruiu praticamente todo o \emph{pipeline}. A deteção passou a assentar num
modelo YOLOv5, e a calibração resulta de duas etapas encadeadas, uma primeira estimativa
dos parâmetros da câmara por um modelo híbrido de rede convolucional e \emph{transformer},
seguida de um refinamento por um modelo de pontos de referência do campo, uma ResNet50 com
74 pontos, arbitrado por um mecanismo de votação sobre o erro de reprojeção. O
\emph{tracking} organiza-se em três fases, a saber, deteção de equipa, seguimento em bruto
e pós-processamento. A função de cada interveniente é decidida por um modelo dedicado de
identificação de equipa, uma OSNet treinada sobre 111 equipamentos distintos, que atribui
uma de cinco funções possíveis. A dimensão do esforço é reveladora, pois o salto para os
63.81\% resulta de intervir em todas as etapas em simultâneo, incluindo o treino de modelos
de raiz sobre dados próprios.

A equipa \textbf{UPCxMobius} (43.15\%), que apresentou o seu trabalho sob o título
\emph{No bells, More Whistles}, seguiu uma estratégia mais próxima da adotada nesta
dissertação, uma vez que partiu da TrackLab e se limitou a substituir três componentes.
Trocou o método de calibração por um processo hierárquico de geração de pontos de
referência que explora as propriedades geométricas do campo, apoiado numa rede
codificador-descodificador HRNetv2. Retreinou o detetor de pessoas sobre cerca de 2000
imagens anotadas manualmente a partir de vídeos do SoccerNet, passando a classificar
diretamente cada deteção como jogador, guarda-redes ou árbitro. Por fim, treinou um
detetor dedicado de dígitos de camisola sobre dados públicos e cerca de 1000 imagens
próprias, combinando depois os dígitos reconhecidos para formar o número. Vale a pena
sublinhar o contraste com a via seguida neste trabalho, pois onde esta equipa optou por
um detetor de dígitos treinado para o efeito, a presente dissertação recorre a um modelo
multimodal genérico guiado pela pose.

A equipa \textbf{Robo Space} (26.92\%) manteve igualmente a TrackLab como estrutura de
base, afinou um YOLOv8 sobre um conjunto próprio, refinou o modelo de reidentificação
BPBReID e acrescentou um mecanismo de atenção temporal global, seguido de suavização
temporal das deteções para eliminar ruído.

A equipa \textbf{playbox x NUSG} (21.26\%) é interessante menos pelo resultado do que
pelo diagnóstico que produziu. Afinou um YOLOv8X para a deteção e corrigiu um defeito na
implementação da calibração, que excluía o círculo central do cálculo do erro de
reprojeção, acrescentando ainda uma filtragem temporal que usa o ângulo de rolamento como
indicador de erro, partindo do princípio de que uma câmara de transmissão não deve
apresentar rolamento. Para o número de camisola construiu um método em duas fases assente
no CLIP, com um modelo a detetar a presença do número e outro a identificá-lo, afinados
sobre 3000 caixas anotadas manualmente, e recorreu ao GPT-4o para refinar as funções, as
equipas e os números. Apesar de uma melhoria de 4\% no conjunto de teste, o sistema não
superou a \emph{baseline} no conjunto de \emph{challenge}, e os próprios autores concluem
que a calibração e a leitura do número continuam a ser as etapas com maior margem.

A equipa \textbf{eidos.ai} (8.92\%) deixou o registo mais instrutivo de todos. Concentrou
o esforço na classificação de função e de equipa, resolvendo-a globalmente para o vídeo em
vez de decidir tracklet a tracklet, através de uma rede siamesa que produz representações
das caixas delimitadoras, de um agrupamento por k-means com vários valores de k e de uma
floresta aleatória que escolhe o agrupamento mais plausível a partir de características de
contexto de jogo. O resultado ficou abaixo da \emph{baseline}, e os autores atribuem-no ao
facto de o modelo ter sido treinado com posições da anotação de referência e avaliado sobre
deteções reais, necessariamente imprecisas. É um exemplo claro de como a qualidade das
etapas anteriores condiciona tudo o que vem depois, e não uma falha do método em si.

\subsection{EDIÇÃO DE 2025}

A segunda edição contou com 14 equipas e 66 submissões. A \emph{baseline}, entretanto
melhorada para 29.01\%, foi superada por 12 equipas, e o melhor resultado subiu ligeiramente
para 63.90\%.

A equipa vencedora, \textbf{KIST-GSR} (63.90\%), deteta os intervenientes com um YOLO-X,
segue-os com o Deep-EIoU apoiado em representações de reidentificação da OSNet e projeta-os
no campo através de um modelo de pontos de referência que considera vários \emph{frames}.
A contribuição mais distintiva está na forma como trata a identidade, formulada como uma
tarefa de geração autorregressiva multimodal resolvida pelo LLaMA-3.2-Vision, que estima a
função, o número de camisola e a cor a partir de instruções em linguagem natural, o que lhe
confere capacidade de reconhecimento em conjunto aberto. Os \emph{tracklets} iniciais são
depois refinados por uma estratégia de divisão e fusão, em que as divisões decorrem da
consistência das identidades previstas e as fusões da semelhança de reidentificação sob
restrições espácio-temporais. A equipa é atribuída por agrupamento da cor da camisola
seguido da comparação da posição horizontal média de cada grupo, assumindo que os jogadores
tendem a manter-se do seu lado, e o resultado é consolidado por votação maioritária dentro
de cada tracklet.

A equipa \textbf{lianyou} (62.76\%) deteta pessoas com um YOLOv12 e as linhas do campo com
um RT-DETR próprio, treinado para 27 categorias de linha. A calibração combina uma correção
de distorção radial com o PnLCalib, reduz os parâmetros da câmara a quatro grandezas e
suaviza-os no tempo. O seguimento parte do IOF-Tracker e melhora a previsão de movimento do
filtro de Kalman com fluxo ótico de Farneback, sobretudo nos movimentos bruscos da câmara.

A equipa \textbf{Playbox \& MIXI} (61.64\%) voltou a partir da TrackLab, substituindo os
módulos por alternativas recentes. Usa o RF-DETR na deteção, que reporta superar a família
YOLO, representações CLIP-ReID na reidentificação e uma cadeia de seguimento composta pelo
BoT-SORT e pela associação global de \emph{tracklets}. Para o número de camisola isola a
região do tronco a partir dos pontos de interesse do ViTPose e reconhece os dígitos com o
PARSeq, depois de rejeitar valores atípicos. A calibração estende o BroadTrack propagando
os parâmetros ao longo do tempo por fluxo ótico. Merece nota a forma como resolve a função
de guarda-redes, atribuída por posição, ao jogador mais afastado dentro das áreas, seguida
de uma otimização que penaliza configurações impossíveis, como a existência de dois
guarda-redes do mesmo lado.

A equipa \textbf{eidos.ai} (46.24\%) apresenta a abordagem mais próxima da adotada nesta
dissertação. Recorre ao Qwen2-VL-Instruct, da mesma família do modelo aqui utilizado, mas
em vez de classificar cada deteção isoladamente entrega ao modelo uma sequência longa de
deteções dispostas em grelha e deixa-o raciocinar sobre elas, com a possibilidade de invocar
uma função de ampliação para inspecionar intervalos específicos antes de decidir o número.
O mesmo mecanismo é usado para identificar o árbitro, e o agrupamento foi alterado para
procurar um melhor equilíbrio entre o número de \emph{tracklets} de cada equipa.

A equipa \textbf{UTokyo Football Lab} (40.13\%) trabalhou sobretudo a qualidade da imagem
de entrada. Aplicou remoção de desfoque com o DeblurGAN-v2, afinou um YOLOv8x, segmentou
as deteções para isolar a silhueta do jogador e dela extraiu um descritor de dez dimensões
com as cores do equipamento e as respetivas áreas, que passou a alimentar a reidentificação,
a função e a equipa. É particularmente relevante para esta dissertação a conclusão com que
os autores encerram, ao reconhecerem que, tal como na \emph{baseline}, a maior dificuldade
permaneceu no reconhecimento do número de camisola, e ao apontarem os modelos multimodais
como o caminho a seguir.

A equipa de \textbf{Hyungjung Kim et al.} (36.90\%) introduziu uma ideia de natureza
diferente, ao explorar uma restrição do próprio domínio. Partindo do princípio de que num
lance não existem mais do que 24 identidades distintas, 22 jogadores e 2 árbitros, usou
essa unicidade para tornar a associação de identidades mais fiável, combinando-a com
remoção de desfoque, compensação do movimento da câmara por fluxo ótico e um detetor
afinado sobre o SoccerNet-GSR.

Por fim, a equipa do \textbf{Laboratory for Biosignal Processing} (28.45\%) dedicou-se
exclusivamente à calibração, refinando a homografia ao longo do vídeo. Identifica os
\emph{frames} de inicialização fiáveis comparando uma segmentação do campo por U-Net com a
que resulta da homografia estimada pelo NBJW, propaga os pontos por fluxo ótico de
Lucas-Kanade e atualiza a homografia por RANSAC, reportando ganhos até 4\% sobre a
calibração de referência, sobretudo nos \emph{frames} pobres em linhas.

\subsection{LEITURA TRANSVERSAL E POSICIONAMENTO DESTE TRABALHO}

Percorridas as duas edições, destacam-se quatro padrões, e todos eles ajudam a situar o
que aqui se propõe.

O primeiro é que o retreino do detetor se tornou prática corrente. Quase todos os sistemas
acima afinam o detetor sobre dados próprios ou anotados de raiz, e os organizadores da
edição de 2025 registam-no explicitamente como norma. Ora, a deteção é precisamente o
módulo que esta dissertação decidiu não tocar, e é razoável admitir que boa parte da
distância que separa os 60\% dos primeiros classificados do valor aqui obtido se explique
por essa diferença de âmbito, e não pela qualidade dos módulos de identificação.

O segundo é o peso da calibração. Duas equipas dedicaram-lhe a totalidade do seu esforço e
quase todas as restantes a substituíram ou reforçaram, com destaque para a introdução do
fluxo ótico como forma de estabilizar os parâmetros ao longo do tempo. Também aqui o
presente trabalho manteve o módulo original.

O terceiro é a adoção generalizada de modelos multimodais de visão e linguagem para a
leitura do número. Em 2024 as abordagens ainda passavam por detetores de dígitos treinados
para o efeito ou por modelos do tipo CLIP, ao passo que em 2025 os organizadores registam o
uso bem-sucedido do LLaMA-Vision e do Qwen2-VL-Instruct, e a equipa vencedora resolve toda a
identidade por essa via. A opção tomada nesta dissertação, decidida de forma independente e
por razões de restrição de memória, coincide portanto com a direção que o estado da arte
veio a seguir.

O quarto padrão é o reconhecimento, por várias equipas, de que a leitura do número continua
a ser o elo mais fraco. A equipa da Universidade de Tóquio afirma-o de forma direta, e a
playbox chega à mesma conclusão pela via da experimentação. Isto sustenta, com fonte
independente, a escolha de âmbito feita neste trabalho.

Resta uma diferença que merece registo. Nenhuma das descrições publicadas destas soluções
menciona um passo de recorte da região dorsal guiado pela estimativa de pose antes de
entregar a imagem ao modelo de leitura. As equipas de 2025 que usam modelos multimodais
alimentam-nos com o recorte do jogador ou com sequências de deteções, e a única que isola
o tronco, a Playbox \& MIXI, fá-lo para alimentar um motor de reconhecimento de texto
convencional e não um modelo multimodal. O contributo aqui proposto situa-se exatamente
nesse ponto, ou seja, na conjugação do recorte anatómico guiado por pose com a leitura
multimodal e a agregação temporal por votação.
